\begin{remark}
\label{remark-tohoku-category-conventions}
\begin{reference}
\cite[Chapter I, \S 1.1, pp. 122--123]{Tohoku}
\end{reference}
\begin{history}
The cited source starts with a nonempty class of objects and prudently
requires the Hom sets for distinct ordered pairs to be disjoint.  The
convention in Definition \ref{definition-category} instead uses a set of
objects, permits the empty category, and treats disjointness of Hom sets as
customary but optional.
\end{history}
These are convention choices rather than a mathematical contradiction.  In
the small nonempty subcase, replacing each arrow by a tagged triple
\((\text{source},\text{target},\text{arrow})\) gives the customary disjoint
presentation of the same category data, so Definition
\ref{definition-category} can be used after this translation.  This remark
does not assert an equivalence for the unrestricted source convention and
does not import its proper-class or nonemptiness requirements into the
Stacks definition.
\end{remark}
